\section{Related work}
\label{sec:related}

\paragraph{Short-term reversal and pairs trading.}
Short-horizon reversal in individual stock returns is well documented: stocks that performed
poorly over the past week or month tend to outperform in the following period
\citep{lehmann1990,lo1990}. \citet{lo1990} show that a large part of contrarian profits comes
from cross-stock effects and not only from own-stock overreaction. This motivates looking at
returns \emph{relative to} related stocks rather than in isolation. \citet{gatev2006} document
the profitability of a simple distance-based pairs-trading rule over several decades of US data.
\citet{krauss2017} surveys the pairs-trading literature and classifies approaches into distance,
cointegration, time-series, stochastic-control and other methods.

\paragraph{Factor models and PCA.}
Arbitrage pricing theory \citep{ross1976} justifies decomposing returns into exposures to a
small number of pervasive factors plus idiosyncratic noise. \citet{connor1986} show that the
factors can be estimated by principal components when the cross-section is large. How many
factors to keep is a long-standing question. \citet{bai2002} propose information criteria for
it, while random-matrix theory offers a simple benchmark: under the null of pure noise, the
eigenvalues of a sample correlation matrix are confined below the Marchenko--Pastur edge
\citep{marchenko1967}. \citet{laloux1999} and \citet{plerou2002} find that only a small number of
eigenvalues of stock correlation matrices exceed this edge. The rest are statistically
indistinguishable from noise.

\paragraph{PCA statistical arbitrage.}
\citet{avellaneda2010} combine these ingredients: PCA (or sector ETFs) for the factors, an OU
model for the residuals \citep{uhlenbeck1930}, and a threshold rule on a standardised residual,
the \emph{s-score}. They also introduce a volume-adjusted (``trading time'') variant. Their
strategy is the baseline replicated here. \citet{khandani2007} analyse the abrupt losses of
quantitative equity market-neutral funds in August 2007, a reminder that such strategies are
crowded and exposed to sudden deleveraging.
\todo{Add post-2010 replications and extensions, e.g.\ deep-learning variants of residual
mean-reversion trading.}

\paragraph{Trading costs and evaluation.}
With proportional costs, the optimal policy is to leave the portfolio unchanged while it stays
inside a no-trade region \citep{davis1990}. \citet{garleanu2013} show that the optimal portfolio
moves only partially towards its target each period. Both ideas motivate the no-trade band in
\cref{sec:method:costs}. On evaluation, \citet{harvey2016} argue that the large number of
strategies tried by researchers calls for much higher significance thresholds. \citet{bailey2014}
propose the deflated Sharpe ratio, which corrects an observed Sharpe ratio for the number of
trials and for non-normal returns. We count every variant we try for exactly this reason.
